\subsection{Uniform independent sentences}
\label{subsec:uniform-independent-sentences}

Fix a presentation $(A,R)$ of a consistent theory with a c.e.\ axiom
set $A$ and a uniformly c.e.\ family $R$ of finitary inference rules.
In Pour-El's terminology \cite{PourEl}, the presentation is
\emph{effectively extensible} if there is a total recursive function $g:\N\to\N$ with the following
property. Whenever $e$ indexes a c.e.\ set $W_e$ of sentences in the
same language and $A\cup W_e$ is consistent under $R$, the number
$g(e)$ codes a sentence $\varphi_e$ such that
$A\cup W_e\nvdash_R\varphi_e$ and $A\cup W_e\nvdash_R\neg\varphi_e$
\cite[pp.~272, 274, Definitions~5, 6(b) and~8]{PourEl1970}.\footnote{The
conclusion of Definition~8 on p.~274 prints $\cap$ for $\cup$. On
p.~272 Pour-El specifies a sentence outside both the theorem set and
the refutation set of the extension.}
For PA we construct this assignment in
$\mathcal B$ using the Turing morphism of $S_{\PA}$.

\begin{theorem}[Uniform independent sentences]
\label{thm:uniform-independent-sentences}
Let $E=(\omega,1_\omega)$ and let $S$ be the arithmetic sentence
object in $\mathcal B=\Tot(\Split(S_{\PA}))$. For each $e\in\N$,
let $\Gamma_e$ consist of the chosen PA axioms and all arithmetic
sentences on whose codes program $e$ halts. There is a morphism
$q:E\to S$ in $\mathcal B$ such that, for every $e\in\N$ with
$\Gamma_e$ consistent, the
sentence with code $(G_{\N}q)(e)$ is neither provable nor
refutable from $\Gamma_e$.
\end{theorem}
\begin{proof}
Use the provability predicate $\operatorname{Prov}(e,z)$ of
Section~\ref{subsec:indexed-arithmetic-theories}, the morphism
$\theta:E\to S$ of Lemma~\ref{lem:arithmetic-diagonal-representation},
and the function $\Theta:\N\to\N$ defined in
Section~\ref{subsec:arithmetic-internal-comparison}.
Define the arithmetic formulas
\[
 C(e,x)\equiv\operatorname{Prov}(e,\Theta(x)),\qquad
 C'(e,x)\equiv\operatorname{Prov}(e,\operatorname{Neg}(\Theta(x))),
\]
and
\[
 R(e,x,z)\equiv
   (C(e,x)\land z=2x+1)\lor(C'(e,x)\land z=2x).
\]
Lemma~\ref{lem:syntactic-functional-search}, with input
$\langle e,x\rangle$, gives a standard index
$a_\ell\in\mathcal I(\omega\times\omega,\omega+\omega)$
such that
\[
 \begin{aligned}
 \PA&\vdash\forall e\,\forall x\,\forall z\,
   \bigl(G_{a_\ell}(\langle e,x\rangle,z)\to R(e,x,z)\bigr),\\
 \PA&\vdash\forall e\,\forall x\,
   \bigl(\exists z\,G_{a_\ell}(\langle e,x\rangle,z)
                     \leftrightarrow C(e,x)\lor C'(e,x)\bigr).
 \end{aligned}
\]
Put $\ell=[a_\ell]_{\PA}:\omega\times\omega\rightharpoonup
\omega+\omega$ in $S_{\PA}$.

Recall the Turing morphism
$\tau_{\PA}:\omega\times\omega\rightharpoonup\omega$,
the isomorphism $b:\omega\to\omega+\omega$, and
$h=b\tau_{\PA}\Delta_\omega:\omega\rightharpoonup\omega+\omega$
in $S_{\PA}$ from Section~\ref{subsec:arithmetic-internal-comparison}.
The Turing property gives a total morphism
$g:\omega\to\omega$ in $S_{\PA}$ satisfying
\[
 \tau_{\PA}(g\times1_\omega)=b^{-1}\ell:
                    \omega\times\omega\rightharpoonup\omega.
\]
Precomposing with $\langle1_\omega,g\rangle:\omega\to\omega\times\omega$
and using $(g\times1_\omega)\langle1_\omega,g\rangle
=\langle g,g\rangle=\Delta_\omega g$ gives
\[
 hg=\ell\langle1_\omega,g\rangle:
                  \omega\rightharpoonup\omega+\omega.
\]
The morphism $I(g):E\to E$ is total, so
$q=\theta I(g):E\to S$ belongs to $\mathcal B$.

Fix $e\in\N$ with $\Gamma_e$ consistent, and put
$n=F_{\N}(g)(e)\in\N$. Choose a first-order structure
$K\models\Gamma_e$ by completeness.
Proposition~\ref{prop:standard-code-soundness} and
Proposition~\ref{prop:arithmetic-theory-fibers} show that
$C(e,x)$ and $C'(e,x)$ have disjoint truth sets in $\N$.
Let $s=F_{\N}(\ell):\N\times\N\rightharpoonup\N\amalg\N$ and
let $i_0,i_1$ be the coproduct injections. For every $x\in\N$ the
graph and domain equations for $a_\ell$, together with disjointness, give
\[
 \begin{aligned}
 ((e,x),i_0(x))\in\operatorname{graph}(s)
     &\quad\Longleftrightarrow\quad \N\models C'(e,x),\\
 ((e,x),i_1(x))\in\operatorname{graph}(s)
     &\quad\Longleftrightarrow\quad \N\models C(e,x).
 \end{aligned}
\]
At $x=n$, the identity $hg=\ell\langle1_\omega,g\rangle$
identifies $F_{\N}(h)(n)$ and $s(e,n)$ as partial values.
If $C(e,n)$ holds, the second equivalence gives $n\in D_1$,
so Lemma~\ref{lem:arithmetic-diagonal-representation} gives
$C'(e,n)$. If $C'(e,n)$ holds, the first gives $n\in D_0$,
so the same lemma gives $C(e,n)$. Disjointness rules out both cases.
For the sentence $\varphi$ with code
$(G_{\N}q)(e)=\Theta(n)$, a finite derivation of $\varphi$ or
$\neg\varphi$ from $\Gamma_e$ would witness $C(e,n)$ or $C'(e,n)$,
respectively. Hence $\varphi$ is independent of $\Gamma_e$.
\end{proof}

